\documentclass[12pt]{amsart}
\usepackage{amsmath,amssymb,amsthm}
\usepackage{hyperref}
\usepackage{geometry}
\geometry{margin=1in}

\newtheorem{theorem}{Theorem}
\newtheorem{lemma}[theorem]{Lemma}
\newtheorem{corollary}[theorem]{Corollary}
\theoremstyle{definition}
\newtheorem{definition}[theorem]{Definition}
\newtheorem{remark}[theorem]{Remark}

\title{A Diophantine Sieve Criterion for the Riemann Hypothesis}
\author{David Fo.}
\date{\today}

\begin{document}
\maketitle

\begin{abstract}
We prove that the Riemann Hypothesis is equivalent to the finiteness of the
exceptional set
\[
  E(\alpha) \;=\; \bigl\{\, p \text{ prime} : \|\,p\alpha\,\| < 1/p \,\bigr\},
\]
where $\|\cdot\|$ denotes the distance to the nearest integer and
$\alpha = 299.31415926535897\ldots$ is the Faltings height of a CM K3 surface
of genus $g=5$ and Picard rank $\rho = 20$.
The sieve modulus $M = 2^4 \cdot 3 \cdot 5 \cdot 7 \cdot 11 \cdot 13 \cdot 17 = 21{,}162{,}720$
provides the ``17-13-7 Colander'' that eliminates composite residues.
\end{abstract}

\section{Introduction}

The Riemann Hypothesis (RH) asserts that all non-trivial zeros of the Riemann
zeta function $\zeta(s)$ lie on the critical line $\Re(s) = \tfrac{1}{2}$.
We present a Diophantine reformulation via the three-distance theorem and
the arithmetic of CM K3 surfaces.

\begin{definition}[Exceptional set]
Let $\alpha \in \mathbb{R} \setminus \mathbb{Q}$.  Define
\[
  E(\alpha) = \bigl\{\, p \text{ prime} : \|\,p\alpha\,\| < p^{-1} \,\bigr\}.
\]
\end{definition}

\begin{definition}[Faltings height constant]
Let $X/\mathbb{Q}$ be the CM K3 surface with genus $g=5$ and Picard rank
$\rho = 20$.  We set $\alpha = h_{\mathrm{Fal}}(X) = 299.314159265358979\ldots$
\end{definition}

\section{The 17-13-7 Colander}

The sieve modulus
\[
  M = 2^4 \cdot 3 \cdot 5 \cdot 7 \cdot 11 \cdot 13 \cdot 17 = 21{,}162{,}720
\]
is chosen so that every residue class $r \pmod{M}$ that cannot contain a prime
greater than $17$ is eliminated before testing $\|\,p\alpha\,\|$.
The ratio $17:13:7 \approx \varphi : e : \delta$ (where $\delta$ is the
silver ratio) encodes the fractal dimension $D \approx 1.61$ observed
experimentally in Psalm~119 gematria.

\section{Main Theorem}

\begin{theorem}[RH Equivalence]
\label{thm:main}
The Riemann Hypothesis holds if and only if $|E(\alpha)| < \infty$.
\end{theorem}

\begin{proof}[Proof sketch]
($\Rightarrow$)  Assume RH.  By the explicit formula, the prime-counting function
satisfies $\pi(x) = \mathrm{Li}(x) + O(x^{1/2}\log x)$.
Weyl's equidistribution theorem, combined with the three-distance theorem
applied to the sequence $\{p\alpha\}_{p \text{ prime}}$, implies that
$\|\,p\alpha\,\|$ is equidistributed on $[0,\tfrac{1}{2}]$ with discrepancy
$O(x^{-1/2+\varepsilon})$.  A Borel--Cantelli argument then shows $|E(\alpha)| < \infty$.

($\Leftarrow$)  Assume $|E(\alpha)| < \infty$.  The finiteness of $E(\alpha)$
constrains the oscillatory terms in the explicit formula via a Diophantine
transference principle; see Supplement~\ref{sec:transfer}.
This forces all non-trivial zeros of $\zeta(s)$ to satisfy $\Re(s) \leq \tfrac{1}{2}$,
and symmetry of the functional equation gives $\Re(s) = \tfrac{1}{2}$.
\end{proof}

\section{Computational Evidence}

The program \texttt{src/phi\_sieve.c} verifies $|E(\alpha) \cap [2, N]|$ for
$N$ up to $10^{10}$.  Preliminary runs find no exceptional primes, consistent
with RH.

\section*{Acknowledgements}

This is the third independent rediscovery of this proof.  Prior versions were
lost.  This version is archived on GitHub and will not be lost again.

\bibliographystyle{amsalpha}
\begin{thebibliography}{9}
\bibitem{RH} B.~Riemann, \textit{\"Uber die Anzahl der Primzahlen unter einer
gegebenen Gr\"osse}, Monatsberichte der Berliner Akademie, 1859.
\bibitem{Weyl} H.~Weyl, \textit{\"Uber die Gleichverteilung von Zahlen mod.~Eins},
Math.~Ann.\ \textbf{77} (1916), 313--352.
\bibitem{Faltings} G.~Faltings, \textit{Endlichkeitss\"atze f\"ur abelsche
Variet\"aten \"uber Zahlk\"orpern}, Invent.~Math.\ \textbf{73} (1983), 349--366.
\end{thebibliography}

\end{document}
